\section{Post-Processing Images (pvPost)}
\codeword{pvPost} is the automated ParaView image generator. It reads the converged solution, builds the standard set of geometry, surface-contour, slice, line-integral-convolution (LIC) and iso-surface images, and writes them to \codeword{postProcessing/images}. It runs head-less through \codeword{pvbatch} and is driven entirely by the \codeword{pvPostSetup} file copied into each case, so no manual ParaView session is required.

\subsection{Inputs and Outputs}
\codeword{pvPost} expects the solution to have been exported to EnSight format. It reads
\begin{lstlisting}
    EnSight/<caseName>.case
\end{lstlisting}
and uses the latest available time step. Iso-surface images additionally read the \codeword{iso*.vtp} files written under \codeword{postProcessing/surfaces/<latestTime>}. All images are written as \codeword{.jpeg} into
\begin{lstlisting}
    postProcessing/images/<variable>_<imageType>/
\end{lstlisting}
one sub-folder per variable and image type. Slice images include the slice normal, position and an index in the file name so they sort in sweep order.

\subsection{Running pvPost}
\codeword{pvPost} is a ParaView script and must be launched with \codeword{pvbatch} (or \codeword{pvpython}), not the system \codeword{python}, so that the \codeword{paraview.simple} modules are available. From inside the case directory:
\begin{lstlisting}
    cd <caseName>
    pvbatch --mpi /path/to/postUtilities/pvPost.py >> log.pvpost
\end{lstlisting}
On the cluster this is wired into the post-processing job through \codeword{clusterDict} and runs automatically after the solve. There is one optional flag:
\begin{itemize}
    \item \codeword{--meshOnly} - generates only the mesh slice images and skips all field images. Use this to inspect the mesh before a solution exists.
\end{itemize}
\begin{lstlisting}
    pvbatch --mpi /path/to/postUtilities/pvPost.py --meshOnly >> log.pvpostMesh
\end{lstlisting}

\subsection{The pvPostSetup File}
\codeword{caseSetup} copies a \codeword{pvPostSetup} file into every case. Edit it to choose which image sets are produced and how the slices are taken. The default file is shown below.
\begin{lstlisting}
    [PV_POST_MAIN]
    CASE_NAME = default
    CAMERA_VIEWS = default
    VARIABLE_DICT = default
    WHEEL_BASE = default
    BACKGROUND_COLOR = default
    RESOLUTION = default
    GEOM_IMG = True
    SURFACE_IMG = True
    SLICE_IMG = True
    SLICE_LIC_IMG = True
    ISO_SURFACE_IMG = True
    CREATE_SCALAR_BAR = True

    [GEOM]
    VIEWS = default

    [SURFACE]
    VIEWS = default
    VARIABLES = default

    [SLICE]
    VIEWS = default
    NORMALS = default
    NSLICES = default
    SLICE_RANGE = default
    VARIABLES = default
    LIC_VARIABLES = default

    [ISO_SURFACE]
    VIEWS = default
    VARIABLES = default
    ISO_VALUES = default

    [VARIABLES]
    SURFACE_VARIABLES = default
    SLICE_VARIABLES = default
    ISO_VARIABLES = default
\end{lstlisting}

\subsubsection{[PV\_POST\_MAIN]}
\begin{itemize}
    \item \codeword{CASE_NAME} [\codeword{str}] - case name used in image titles and file names, \codeword{default} uses the case directory name (\codeword{default}, or \codeword{str})

    \item \codeword{CAMERA_VIEWS} [\codeword{str}] - path to a camera-views file, \codeword{default} uses the bundled view definitions (\codeword{default}, or a file path)

    \item \codeword{VARIABLE_DICT} [\codeword{str}] - path to the variable dictionary that defines the plotted variables, their equations, labels, ranges and colour maps, \codeword{default} uses \codeword{postUtilities/default/defaultVariables} (\codeword{default}, or a file path)

    \item \codeword{RESOLUTION} [\codeword{str}] - working render resolution, \codeword{default} uses the built-in 16:9 setting (\codeword{default})

    \item \codeword{GEOM_IMG}, \codeword{SURFACE_IMG}, \codeword{SLICE_IMG}, \codeword{SLICE_LIC_IMG}, \codeword{ISO_SURFACE_IMG} [\codeword{bool}] - master switches for each image set, set to \codeword{False} to skip that set (\codeword{True}, \codeword{False})

    \item \codeword{CREATE_SCALAR_BAR} [\codeword{bool}] - draw the colour-bar legend on field images (\codeword{True}, \codeword{False})
\end{itemize}

\subsubsection{[SLICE]}
The slice block controls both the field slices and the LIC slices. Each of \codeword{VIEWS}, \codeword{NORMALS}, \codeword{NSLICES} and \codeword{SLICE_RANGE} is a list, and the lists must all be the same length: entry $i$ defines one slice plane sweep.
\begin{itemize}
    \item \codeword{VIEWS} [\codeword{str}] - camera view used to render each plane (\codeword{default} = \codeword{Front Left LeftForward LeftBack Top})

    \item \codeword{NORMALS} [\codeword{str}] - slice plane normal for each sweep (\codeword{default} = \codeword{X Y Y Y Z})

    \item \codeword{NSLICES} [\codeword{int}] - number of evenly spaced slices in each sweep, \codeword{default} adapts to the symmetry of the case

    \item \codeword{SLICE_RANGE} [[\codeword{float,m}] [\codeword{float,m}]] - start and end coordinate of each sweep as \codeword{[min,max]}, \codeword{default} derives sensible ranges from the front-axle reference and the reference dimensions

    \item \codeword{VARIABLES} [\codeword{str}] - field-slice variables to plot, \codeword{default} plots the standard pressure, velocity and vorticity set

    \item \codeword{LIC_VARIABLES} [\codeword{str}] - variables coloured underneath the LIC streamlines, \codeword{default} uses the standard set
\end{itemize}
A variable requested here that is not defined in the variable dictionary, or not present in the solution, is reported and skipped rather than stopping the run.

\subsubsection{[SURFACE] and [ISO\_SURFACE]}
\begin{itemize}
    \item \codeword{[SURFACE]} \codeword{VIEWS} / \codeword{VARIABLES} - camera views and surface-contour variables painted on the car body (e.g. \codeword{CpMean}, wall shear).

    \item \codeword{[ISO_SURFACE]} \codeword{VIEWS} / \codeword{VARIABLES} / \codeword{ISO_VALUES} - camera views for the iso-surface renders. The iso-surfaces themselves (\codeword{isoQ}, \codeword{isoCtp}, ...) are produced by the solver's surface function objects and read from \codeword{postProcessing/surfaces}; \codeword{pvPost} renders whatever \codeword{iso*.vtp} files it finds.
\end{itemize}

\subsection{Variables and Camera Views}
The variable dictionary defines every plottable variable as an equation over the OpenFOAM fields, together with a display label, a value range and a ParaView colour preset. The tokens \codeword{UREF}, \codeword{LREF} and \codeword{WREF} in the equations are substituted with the case reference values before evaluation, so coefficients such as \codeword{CpMean} and \codeword{CptMean} are normalised consistently with the rest of the case. To use a custom set, point \codeword{VARIABLE_DICT} (and, if needed, \codeword{CAMERA_VIEWS}) at your own file.

To customise one case, copy the default variable and view files from \codeword{postUtilities/default/} to new files, then reference those files from the case's \codeword{pvPostSetup}. Do not modify the repository defaults unless the change should apply to every future case. For example:
\begin{lstlisting}
    [PV_POST_MAIN]
    VARIABLE_DICT = myVariables
    CAMERA_VIEWS = myViews

    [SURFACE]
    VIEWS = DriverSide FrontLeft
    VARIABLES = CpMean UMeanStreamwise CfMean

    [SLICE]
    VIEWS = DriverSide
    VARIABLES = CpMean UMeanStreamwise
    LIC_VARIABLES = UMean
\end{lstlisting}

Each custom variable file contains the three top-level groups \codeword{surfaceVariables}, \codeword{sliceVariables} and \codeword{isoVariables}. Add a variable to the required groups using an equation, label, display range and colour map. The following entry adds a normalised streamwise mean velocity variable:
\begin{lstlisting}
    "UMeanStreamwise": {
        "equation": "UMean_X/UREF",
        "label": "Umean,x/Uref",
        "range": [0, 1.2],
        "color": "Viridis (matplotlib)"
    }
\end{lstlisting}
The name \codeword{UMeanStreamwise} is used in the \codeword{VARIABLES} lists. The equation must reference fields present in the exported EnSight solution. A requested variable that is missing from the dictionary or solution is reported and skipped.

Custom camera-view files use one JSON object per view. The camera position is \codeword{campos}, the point being viewed is \codeword{focalpos}, \codeword{viewup} defines the upward direction, and \codeword{parallelscale} controls the orthographic zoom. Reference tokens can be used in the expressions:
\begin{lstlisting}
    "DriverSide": {
        "pp": 1,
        "campos": "[0,-LREF*10,0]",
        "focalpos": "[1.18,0,0.85]",
        "viewup": "[0,0,1]",
        "parallelscale": "WREF/2"
    }
\end{lstlisting}
The name in each \codeword{VIEWS} list must exactly match a key in the custom view file. When \codeword{CAMERA_VIEWS = default}, built-in views are generated automatically from the case reference dimensions.

\subsection{Front- and Rear-Wing Pressure Sections}
\codeword{pvPost} can extract the mean pressure coefficient where y-normal planes intersect the front or rear wing surface. Enable either section in the case's \codeword{pvPostSetup} file:
\begin{lstlisting}
    [FRONT_WING_PRESSURE]
    ENABLE = True
    PATCH_PATTERN = frontWing.*
    NORMAL = default
    Y_RANGE = default
    N_PLANES = 11
    VARIABLE = CpMean
    OUTPUT_DIR = postProcessing/wingPressure

    [REAR_WING_PRESSURE]
    ENABLE = True
    PATCH_PATTERN = rearWing.*
    NORMAL = default
    Y_RANGE = default
    N_PLANES = 11
    VARIABLE = CpMean
    OUTPUT_DIR = postProcessing/wingPressure
\end{lstlisting}
The default patch patterns are case-insensitive regular expressions. The default \codeword{NORMAL} is the SAE y direction, \codeword{(0 1 0)}. \codeword{Y_RANGE = default} samples from \codeword{CREF_y - 0.5 WREF} to \codeword{CREF_y + 0.5 WREF}; an explicit range can be supplied as two values, for example \codeword{[-0.4,0.4]}. \codeword{N_PLANES} controls the number of evenly spaced planes. A different three-component plane normal can be supplied instead of \codeword{default}.

Each intersection is written as a CSV file containing \codeword{x}, \codeword{y}, \codeword{z} and \codeword{CpMean} columns under \codeword{OUTPUT_DIR}. Both sections are disabled by default. This feature writes CSV data only and does not alter the existing surface, slice or LIC image generation.

For a generated cluster case, enable the required sections in the case copy of \codeword{pvPostSetup} before submitting the post-processing job. The standard \codeword{postPro} workflow runs \codeword{pvPost.py} first, creating the CSV files when the sections are enabled, and then runs \codeword{postRun.py --wingPressure} to create the plots. No additional cluster setting is required beyond enabling the relevant section; the command is included in the \codeword{post} section of the selected \codeword{clusterDict} template.

For local or manual processing, run \codeword{postRun.py} after \codeword{pvPost.py} has completed. The selected cases are supplied with \codeword{--trial}:
\begin{lstlisting}
    python postUtilities/postRun.py --wingPressure \
        --trial case001 case002 --saveFormat png
\end{lstlisting}
For each selected case, \codeword{postRun.py} reads the available wing CSV files, plots \codeword{CpMean} against \codeword{x} for each y plane, and writes \codeword{frontWing_CpMean_vs_x.png} and \codeword{rearWing_CpMean_vs_x.png} under \codeword{postProcessing/wingPressure}. The pressure-coefficient axis is inverted using the conventional aerodynamic presentation. Cases without enabled wing extraction or matching CSV files are reported and skipped.

\subsection{Half-Symmetry Cases}
When the case is built on half symmetry (\codeword{SIM_SYM}\,=\,\codeword{half}), \codeword{pvPost} automatically reflects the body surface, the internal volume and the iso-surfaces about the symmetry plane so the rendered images show the full car.

\subsection{Cornering (SRF) Cases}
For an SRF cornering case the time-averaged velocity stored in the solution is the rotating-frame relative velocity \codeword{UrelMean}, not the absolute \codeword{UMean}. \codeword{pvPost} detects a cornering case automatically from \codeword{CORNERING_SETUP} -> \codeword{RUN_CORNERING} and, when it is \codeword{True}:
\begin{itemize}
    \item rewrites every variable equation that references \codeword{UMean} (including its components \codeword{UMean_X/Y/Z} and the \codeword{UMean^2} term in \codeword{CptMean}) to read \codeword{UrelMean} instead, and

    \item uses \codeword{UrelMean} as the vector field for the LIC streamlines.
\end{itemize}
Variable names, labels and ranges are unchanged, so the same \codeword{pvPostSetup} works for both straight-line and cornering runs and the image file names stay consistent. A line is printed to \codeword{log.pvpost} confirming the substitution:
\begin{lstlisting}
    Cornering (SRF) case detected: using UrelMean in place of UMean.
\end{lstlisting}
\textbf{NOTE:} This requires \codeword{UrelMean} to be present in the solution. The cornering setup averages both \codeword{U} and \codeword{Urel}, so \codeword{UrelMean} is available; if it is missing, those variables are reported as unavailable and skipped.

\subsection{Performance}
For speed, \codeword{pvPost} loads only the fields actually referenced by the active variable equations (plus the LIC vector) from the EnSight reader, dropping unused fields such as \codeword{phi}, \codeword{k}, \codeword{omega}, \codeword{nut} and the \codeword{Prime2Mean} tensors before the cell-to-point conversion and slicing. It also reuses a single slice filter and calculator per slice normal, updating only the slice position as it sweeps, rather than rebuilding the pipeline for every slice. These changes reduce runtime and memory with no change to the images produced. The console reports which fields were loaded and which were skipped:
\begin{lstlisting}
    Loading only required fields: ['pMean', 'UMean', ...]
    Skipping unused fields: ['phi', 'k', 'omega', ...]
\end{lstlisting}

\subsection{Workflow Summary}
\begin{itemize}
    \item Export the converged solution to EnSight (\codeword{EnSight/<caseName>.case}).

    \item Edit \codeword{pvPostSetup} to choose the image sets (\codeword{GEOM_IMG}, \codeword{SURFACE_IMG}, \codeword{SLICE_IMG}, \codeword{SLICE_LIC_IMG}, \codeword{ISO_SURFACE_IMG}) and, if needed, the slice \codeword{NORMALS}, \codeword{NSLICES} and \codeword{SLICE_RANGE}.

    \item Run \codeword{pvPost} with \codeword{pvbatch} from the case directory (automatic on the cluster). Add \codeword{--meshOnly} to render only the mesh.

    \item Collect the images from \codeword{postProcessing/images}. Cornering cases are handled automatically; no setup change is required.
\end{itemize}

\subsection{Troubleshooting}
\begin{itemize}
    \item \codeword{ERROR! Unable to find EnSight File!} - no \codeword{EnSight/<caseName>.case} was found. Export the solution to EnSight before running \codeword{pvPost}.

    \item \codeword{The inputs in pvPostSetup is not correct for [SLICE]} - the \codeword{VIEWS}, \codeword{NORMALS}, \codeword{NSLICES} and \codeword{SLICE_RANGE} lists are not all the same length. Make every \codeword{[SLICE]} list have one entry per sweep.

    \item \codeword{<variable> not available, skipping} - the variable is not defined in the variable dictionary or not present in the solution. Check the spelling in \codeword{pvPostSetup} and confirm the field was written by the solver.

    \item \codeword{Could not determine required fields from variables, loading all} - field detection found nothing to load and fell back to loading every field. The run still completes; check that the variable dictionary equations reference real field names.

    \item \codeword{No iso-surface files found, skipping} - no \codeword{iso*.vtp} files exist under \codeword{postProcessing/surfaces}. Enable the iso-surface function objects in the solver setup if iso images are required.
\end{itemize}

\section{Case Summaries (postRun.py --summary)}
\codeword{postRun.py --summary} reads the converged solver logs and function-object output for a case and writes a single \codeword{summary.csv} file at the case root. \codeword{summary.csv} is a two column, header-less file of \codeword{key,value} pairs, one row per reported quantity (force and moment coefficients, cell count, mesher, solve time, run date, and so on). Every numeric value is formatted to three decimal places before being written, so downstream tools and spreadsheets show a consistent precision regardless of the underlying solver output.

\begin{lstlisting}
    python postUtilities/postRun.py --summary
\end{lstlisting}

\subsection{Case Completeness}
Before summarising, \codeword{postRun.py} checks that the case actually finished by looking for a standalone \codeword{End} line in the relevant solver log. It checks, in order, \codeword{log.simpleFoam}, \codeword{log.pisoFoam}, \codeword{log.SRFSimpleFoam} and \codeword{log.SRFPimpleFoam}, so both straight-line and cornering (SRF) cases are recognised as complete. A case without a matching, finished log is reported as incomplete and skipped wherever a completeness check is required (single-case summaries, ride-height parent averaging and sensitivity plotting).

\subsection{Ride-Height Parent Summaries}
When \codeword{--summary} is run inside a ride-height mapping parent case (a directory containing child cases named \codeword{<parentCaseName>\_<number>}), \codeword{postRun.py} does not recompute coefficients from the raw logs. Instead it reads each child's own \codeword{summary.csv}, skips any child that is incomplete, missing its \codeword{summary.csv}, or has an unreadable \codeword{summary.csv} (with a warning printed for each skipped child), and averages the remaining numeric fields into the parent's own \codeword{summary.csv}. This includes the lift-balance and side-force coefficients \codeword{Cl(f)}, \codeword{Cl(r)}, \codeword{Cs(f)} and \codeword{Cs(r)} alongside \codeword{Cd} and \codeword{Cl}.

\section{Ride-Height Sensitivity Plots (postRun.py --sensitivityPlots)}
For a ride-height mapping parent case, \codeword{postRun.py --sensitivityPlots} detects and plots single-variable sensitivity sweeps directly from the parent's \codeword{rideHeights\_updated.csv} and each child's \codeword{summary.csv}. It only runs for parent cases (those with child directories matching \codeword{caseName\_\#}); it does nothing for a plain single case or when run from inside a child case.

\begin{lstlisting}
    python postUtilities/postRun.py --sensitivityPlots
    python postUtilities/postRun.py --sensitivityPlots --includeSideForce
    python postUtilities/postRun.py --sensitivityPlots \
        --sensitivityCases ../otherRideHeightMap ../anotherRideHeightMap
\end{lstlisting}

\subsection{Sweep Detection}
The ride-height map columns are organised into groups: \codeword{Front Ride Height} (\codeword{fl}, \codeword{fr}), \codeword{Rear Ride Height} (\codeword{rl}, \codeword{rr}), \codeword{Yaw} (\codeword{yaw}) and \codeword{Corner} (\codeword{corner\_radius}, \codeword{corner\_dir}). The \codeword{steer} column (or \codeword{steer\_deg}/\codeword{steer\_angle}) is tracked separately as a constancy-only group: it is never itself plotted, but it must stay constant for a group to qualify as a clean single-variable sweep. Within a group, columns that move together (for example \codeword{fl} and \codeword{fr} changing in step) are averaged into one representative value, so both changing together still counts as a single variable rather than two.

Because a ride-height map is often a full grid (front ride height crossed with rear ride height, for example), a group can produce several sweeps at different fixed values of everything else. \codeword{postRun.py} finds every fixed-combo subset of the child cases where all other groups (including \codeword{steer}) hold one constant value, and where the named group still varies across at least three of those rows; each qualifying subset is plotted as its own sweep. For example, a map that varies both front and rear ride height produces one Front Ride Height sweep for each rear ride height value that has enough points, and vice versa.

\subsection{Sweep Plots}
Each sweep is drawn as one figure with one subplot per force coefficient: \codeword{Cd}, \codeword{Cl}, \codeword{Cl(f)} and \codeword{Cl(r)} by default, with \codeword{Cs(f)} and \codeword{Cs(r)} added when \codeword{--includeSideForce} is given. Each subplot shows the raw data points plus a low-order polynomial curve fit (degree scales with the number of available points, capped at three, and falls back to a straight line when there are too few unique x-values to fit safely). The figure title names only the swept variable; the values held constant for that sweep (the other groups and \codeword{steer}) are listed in a boxed side panel instead, so long configuration lists do not overrun the plot. Each figure is saved to
\begin{lstlisting}
    postProcessing/sensitivityPlots/<SweepName>_sweep.png
\end{lstlisting}
with a suffix built from the fixed configuration values when more than one sweep of the same type exists (for example \codeword{FrontRideHeight\_rl-0.01\_sweep.png}).

\subsection{Comparing Multiple Ride-Height Maps}
The optional \codeword{--sensitivityCases} flag takes one or more other ride-height mapping parent case paths. For each of those cases, \codeword{postRun.py} builds its own sweeps in the same way, then overlays a sweep from another case onto the primary case's matching sweep only when both the sweep type and every fixed configuration value are identical (same held rear ride height, yaw, corner and steer, for example). Matching sweeps are drawn in the same figure with one colour per case and a legend; sweeps that do not have a matching fixed configuration in another case are left as single-case plots. This makes it possible to compare, for example, a Front Ride Height sensitivity between two different aero concepts at the same rear ride height.

\subsection{Front-Rear Ride Height Contour Plots}
When the map contains a rectangular grid of front and rear ride height values with roll held at zero, \codeword{postRun.py} additionally produces a Front-Ride-Height-vs-Rear-Ride-Height contour plot for each qualifying fixed combination of the remaining variables (Yaw, Corner and steer). The contour requires at least four grid points with at least two unique front and two unique rear values. Each force coefficient is interpolated over the front/rear grid using cubic interpolation (\codeword{scipy.interpolate.griddata}, \codeword{method='cubic'}) and rendered as a filled contour with contour lines and the underlying data points overlaid, arranged two-by-two across the default four metrics (six when \codeword{--includeSideForce} is set, with any unused subplot slots hidden). Output is saved to
\begin{lstlisting}
    postProcessing/sensitivityPlots/FrontRearRideHeight_contour<suffix>.png
\end{lstlisting}
Contour plots are always single-case; they are not affected by \codeword{--sensitivityCases}.

\section{Pushing Results to Google Sheets (gsheetExport.py)}
\codeword{postUtilities/gsheetExport.py} pushes a case's summary metrics to a Google Sheet, one row per trial, using a service account for authentication. It reads its values directly from the case's \codeword{summary.csv} rather than recomputing anything from the solver output, so it must be run after \codeword{postRun.py --summary} has produced an up to date \codeword{summary.csv}.

When the case being pushed is itself a child of a ride-height mapping parent (its directory name matches \codeword{<parentCaseName>\_<number>}), \codeword{gsheetExport.py} first regenerates the parent's \codeword{summary.csv} by invoking \codeword{postRun.py --summary} in the parent directory, then pushes both the child's row and the parent's (map-averaged) row to the sheet. For a case that is not part of a ride-height map, only that case's own row is pushed. The fixed set of spreadsheet columns written for each row includes run date, solve time, cell count, mesher, half/full symmetry, velocity, yaw, density, reference area and wheelbase, \codeword{Cd}, \codeword{Cl}, \codeword{Cl(f)}, \codeword{Cl(r)}, \codeword{Cs(f)}, \codeword{Cs(r)}, confidence intervals for \codeword{Cd}/\codeword{Cl}, and the front and rear wing \codeword{Cd}/\codeword{Cl}.
